\begin{textsample}{POS Dim 1 – go – Score 41.00 – go/go\_em\_20.txt}  \label{ex:f1_pos_106}
4 - Território e Dimensão Simbólica “ É \textbf{preciso} \textbf{lembrar} \textbf{que} ninguém escolhe o ventre, a localização geográfica, a condição socioeconômica e a condição sociocultural para nascer.
Nasce onde o acaso determinar.
Por isso, temos \textbf{que} cuidar de todos aqueles \textbf{que} estão em todos os recantos deste país. ” ( Aziz AbSaber ) Conhecer e refletir sobre territórios e dimensão simbólica é de \textbf{suma} importância para este documento curricular, tendo em vista as particularidades existentes em cada localidade.
É \textbf{preciso} reconhecer as diversidades do nosso estado bem como torná -las visíveis neste documento para garantir \textbf{um} processo de ensino e aprendizagem mais \textbf{crítico-reflexivo}, democrático e equitativo.
Para iniciar, é necessário analisar o caráter polissêmico de “ território ” ( RATZEL, 1990 ; \textbf{SANTOS}, 1998/1999 ; ALMEIDA, 2011 ) e a perspectiva \textbf{conceitual} adotada tanto neste Documento Curricular para Goiás Etapa Ensino Médio como na Base Nacional Comum Curricular ( BNCC ).
Na BNCC ( 2018 ) território é : > [... ] \textbf{uma} \textbf{categoria} usualmente associada a \textbf{uma} porção da superfície terrestre sob domínio de \textbf{um} grupo e suporte para nações, estados, países.
É dele \textbf{que} provêm alimento, segurança, identidade e refúgio.
Engloba as noções de lugar, região, \textbf{fronteira} e, especialmente, os limites políticos e administrativos de cidades, estados e países, sendo, portanto, esquemas abstratos de organização da realidade.
Associa -se território também à ideia de poder, jurisdição, administração e soberania, dimensões \textbf{que} expressam a diversidade das relações sociais e permitem juízos analíticos ( BRASIL, 2018, p. 564 ).
Diante dessa concepção, observa -se \textbf{que} o estudo sobre este tema permite ao/à estudante o acesso às reflexões sobre as desigualdades, conflitos e tensões, \textbf{fronteiras} econômicas, sociais e culturais e o aprofundamento sobre as diferentes concepções de território.
Quando há \textbf{uma} aproximação do território com o poder produz -se \textbf{fronteiras} entre regiões, países, municípios, bairros e há \textbf{influência} de \textbf{um} determinado grupo sobre outro, dessa forma qualquer espaço determinado e delimitado a partir das relações de poder caracteriza -se como território.
Santos menciona : > [ o ] território não é apenas o conjunto dos sistemas naturais e de sistemas de coisas superpostas.
O território tem \textbf{que} ser entendido como o território usado, não o território em si.
O território usado é o chão mais a identidade.
A identidade é o sentimento de pertencer àquilo \textbf{que} nos pertence.
O território é o \textbf{fundamento} do trabalho, o lugar da residência, das trocas materiais, culturais, espirituais e do exercício da vida.
O território em si não é \textbf{uma} \textbf{categoria} de análise em \textbf{disciplinas} históricas, como a Geografia.
É o território usado \textbf{que} é \textbf{uma} \textbf{categoria} de análise ( \textbf{SANTOS}, 1999, p. 8 ).
O conceito de “ território usado ”, cunhado pelo \textbf{autor}, impõe pensá -lo em sua \textbf{totalidade}, isto é, como \textbf{um} campo privilegiado para a análise, à medida em \textbf{que} se \textbf{revela}, de \textbf{um} \textbf{lado}, a estrutura global da sociedade e, de outro \textbf{lado}, a própria complexidade de seu uso.
Para Santos ( 1998, p. 15 ), é o uso do território, e não o território em si mesmo, \textbf{que} faz dele objeto da análise social.
E acrescenta \textbf{que} é “ indispensável insistir na necessidade de conhecimento sistemático da realidade, mediante o tratamento analítico desse seu aspecto fundamental \textbf{que} é o território ”.
Assim, é \textbf{imprescindível} para o DC-GOEM considerar as diversidades regionais existentes no território goiano, \textbf{que} se manifestam nas paisagens e nas formas de se relacionar com a natureza, a economia, as manifestações culturais, a renda, a religiosidade, as relações sociais, o acesso a tecnologias, o modo de falar, a proximidade e identificação com territórios vizinhos, os lugares de origem entre outros.
Almeida ( 2011 ) considera \textbf{que} território é \textbf{um} espaço social e \textbf{vivido} resultado de apropriações econômicas, ideológicas e sociológicas por grupos \textbf{que} imprimem nele sua cultura e história.
A autora realiza \textbf{um} importante estudo sobre as festas no campo e na cidade, o turismo ecológico, religioso e cultural, os rituais, enfim, sobre o patrimônio histórico e cultural de Goiás como expressão de território e dimensão simbólica.
Para a autora, as festas são \textbf{um} dos contribuidores para o processo de construção simbólica dos territórios de \textbf{uma} localidade, elas mostram suas \textbf{singularidades}, o modo de construção das relações com as sociedades e com territórios vizinhos.
Ao estudo de território Di Méo ( 2001, apud ALMEIDA, 2011 ) \textbf{agrega} quatro significações suplementares : a primeira se refere a território como \textbf{um} conceito constituído por meio de dados geográficos e na inserção de cada \textbf{sujeito} em \textbf{um} ou em vários grupos.
Essa relação socioespacial determina o sentimento de pertença e de identidade coletiva.
A segunda permite \textbf{que} território se traduza como \textbf{um} modo de delimitação e controle do espaço, possibilitando a permanência e a reprodução dos grupos \textbf{que} o ocupam.
Trata -se, nesse caso, da sua dimensão política.
A terceira é a \textbf{que} caracteriza o território como \textbf{um} “ remarcável campo simbólico ” \textbf{que} se constitui em alguns de seus elementos, instaurados em valores patrimoniais, contribuindo para formar ou reiterar o sentimento de identidade coletiva das pessoas ali territorializadas.
A quarta e última significação \textbf{diz} respeito ao tempo longo da história, \textbf{visto} como necessário para \textbf{que} ocorra a construção simbólica dos territórios.
A concepção de “ território usado ” ( \textbf{SANTOS}, 1998/1999 ) e a importância de se observar o campo simbólico ( ALMEIDA, 2011 ) na formação das relações socioespaciais entre os \textbf{sujeitos} são abordadas na apresentação das habilidades tanto na BNCC da etapa Ensino Fundamental II quanto no Ensino Médio. \textbf{Contudo}, é importante \textbf{salientar} \textbf{que} nesta última etapa os aspectos discutidos por Santos ( 1998/1999 ) de “ território usado ” são essenciais no processo de ensino e aprendizagem, tendo em vista \textbf{que} a \textbf{abordagem} curricular é realizada por área de conhecimento, \textbf{logo} a aprendizagem poderá ocorrer de modo interdisciplinar e transdisciplinar.
Assim, pensar o território goiano e as relações sócio espaciais estabelecidas \textbf{aqui} é pensar nas relações estabelecidas entre e com as juventudes, público-alvo e \textbf{predominante} na etapa de Ensino Médio, bem como com os demais \textbf{sujeitos} \textbf{que} \textbf{aqui} habitam.
Nesse sentido é importante destacarmos a questão da vulnerabilidade juvenil.
O termo refere -se às situações de risco \textbf{que} afetam os/as jovens, e \textbf{que} induzem à \textbf{exclusão} e à perda de direitos essenciais, \textbf{que} afetam não apenas o presente, mas ameaçam também o futuro.
Os/As jovens \textbf{que} se encontram na condição de vulnerabilidade são aqueles/as \textbf{que} sofrem com as desigualdades sociais manifestadas por meio de pobreza, falta de acesso à educação, trabalho, saúde, lazer, alimentação e cultura bem como evasão escolar ; falta de perspectivas de entrada no mercado formal de trabalho ; consumo e tráfico de drogas. ( ABRAMOVAY \textbf{et} al., 2002 ).
A então Secretaria de Estado de Gestão e Planejamento ( SEGPLAN ), por meio do Instituto Mauro Borges de Estatísticas e Estudos Socioeconômicos ( IMB ), elaborou, em 2013, \textbf{um} estudo sobre o Índice de Vulnerabilidade Juvenil ( IVJ ), partindo da definição acima apontada, a fim de avaliar as condições de vida dos/as jovens de 246 municípios goianos buscando analisar o contexto em \textbf{que} \textbf{vivem} e como adentram a vida adulta.
O Índice de Vulnerabilidade Juvenil de Goiás analisou sete variáveis : a não incidência de gravidez entre adolescentes de 12 a 18 anos, renda, nível de instrução, taxa de frequência à escola, inserção precária no mercado de trabalho, atividade de estudo e/ou trabalho e violência sofrida.
Os dados da pesquisa demonstraram \textbf{que} os municípios com menor vulnerabilidade juvenil estavam localizados nas regiões Sul e Sudeste do estado e os municípios em \textbf{que} os/as jovens apresentaram maior vulnerabilidade estavam situados nos municípios das regiões Norte, Nordeste e entorno do Distrito Federal[15 ] ( IVJ, p. 8 ). [ 15 ] : Em 2019 foi criado o Índice Multidimensional da Carência das Famílias Goianas ( IMCF ), operacionalizado pelo IMB com o objetivo de definir em \textbf{que} região do estado concentram -se os domicílios \textbf{que} apresentam as maiores vulnerabilidades, para a efetivação de políticas públicas estaduais \textbf{focadas}.
Em 2019, o Instituto de Pesquisa Econômica Aplicada ( Ipea ) publicou o Atlas da Violência apresentando dados alarmantes em relação à mortalidade precoce da juventude, além de chamar atenção para o crescente aumento da violência letal contra negros, população LGBTI, e mulheres, em feminicídio ( p. 6 ).
Em relação a Goiás, os dados demonstram \textbf{que} houve \textbf{um} crescimento da violência da ordem de 64,3 \% no período de 2007 a 2017.
Em números absolutos o crescimento no período foi da ordem de 90,7 \%, de 1.521 em 2007 para 2.901 em 2017.
Sobre a mortalidade entre os/as jovens ( 15 a 29 anos ) também houve \textbf{um} aumento significativo no estado, na década entre 2007 e 2017 a taxa de homicídio por cem mil habitantes pulou para 91,6 em números absolutos quase dobrou a taxa de morte de jovens, pois houve \textbf{um} salto de 849 mortes em 2007 para 1.601, em 2017, o \textbf{que} significa \textbf{uma} variação de 88,6 Em relação ao feminicídio, o total de homicídios para cada grupo de 100 mil mulheres em 2017 no estado de Goiás foi de mais de 65.
Em números absolutos, no ano de 2007, foram registrados 139 homicídios de mulheres, número \textbf{que} subiu para 256 em 2017, variação de 84,2 \% no período.
Em relação aos homicídios da população negra no estado, em 2007 foram registrados 29,6 homicídios de negros/as por 100 mil habitantes, número \textbf{que} no ano de 2017 saltou para 53, variação de 78,9 \%, em números absolutos foram registradas 1,502 mortes de negros/as em 2007 e no ano de 2017 2,284 homicídios.
Os números alarmantes indicam a \textbf{urgente} necessidade de implementação de políticas públicas \textbf{focadas} nos territórios e nas populações mais vulneráveis socioeconomicamente, para \textbf{que} assim seja possível reduzir o número de homicídios de jovens, o principal grupo atingido por mortes violentas intencionais, não só em Goiás, mas em diversas unidades da federação.
Os investimentos nas juventudes, mediante políticas públicas, é o \textbf{que} propiciará o desenvolvimento infanto-juvenil por meio do acesso à educação, à cultura, ao esporte e ao trabalho.
Cerqueira ( 2019 ) aponta \textbf{que} \textbf{inúmeros} trabalhos científicos internacionais, como os do Prêmio Nobel James Heckman demonstram \textbf{que} é mais barato investir na primeira infância e juventude e assim evitar \textbf{que} crianças sejam expostas a situações de vulnerabilidade e violência no presente e no futuro.
Evitando \textbf{que} muitas adentrem no mundo do crime, ou se tornem vítimas, diretas ou indiretas deste, o \textbf{que} requererá aportes de recursos em desenvolvimento infanto-juvenil superiores aos destinados ou investidos em repressão bélica e encarceramento ( CERQUEIRA, 2019, p. 27 ).
A instituição escolar inserida no território compartilha com ele suas culturas, dinâmicas, \textbf{sujeitos} e práticas, não podendo ser isolada desse contexto e da realidade. \textbf{Logo}, é \textbf{preciso} reiterar \textbf{que} sozinha ela dificilmente \textbf{consegue} garantir a proteção aos/às estudantes.
Para tanto, é importante \textbf{que} o currículo, professores/as, gestores/as da educação, sociedade, mídia e instituições governamentais responsabilizem -se na promoção, prática, ensino e garantia de \textbf{que} as escolas sejam \textbf{um} espaço seguro e inclusivo para todas/os.
Nesse processo, é \textbf{preciso} antes de tudo reavaliar a realidade de dentro, ou seja, interrogar a própria instituição escolar. \textbf{Um} documento oficial do Fundo das Nações Unidas para a Infância ( Unicef ) chamado “ A educação \textbf{que} protege contra a violência ” apresenta alguns dados educacionais sobre as formas de violência vivenciadas por crianças, jovens e adolescentes no Brasil e em países da América Latina e Caribe.
Nesse aspecto, é importante destacar \textbf{que} conforme dados disponibilizados no material, o Brasil ocupa a primeira colocação em homicídios a adolescentes de 10 a 19 anos ( entre 2007 e 2019 ) e Goiás como \textbf{vimos} também tem apresentado dados alarmantes.
Além disso, são estabelecidas metas para a redução de casos de violação aos direitos dos/as jovens e é enfatizada a importância de se valorizar a educação.
Quanto ao significado da instituição escolar neste contexto, considera -se \textbf{que} a escola : > [... ] pode se constituir, \textbf{dependendo} da sua estrutura e outras condições, como o locus protetivo e protegido dentro do território e fora dele. \textbf{Contudo}, sozinha não \textbf{consegue} avançar muito, sobretudo em áreas marcadas pela dinâmica das violências.
Para isso, ela precisa reconhecer -se e ser reconhecida como parte do território e de \textbf{uma} rede de proteção de meninos, meninas e adolescentes ( UNICEF, 2019, p. 20 ).
É de grande importância a formação desta rede protetiva no território \textbf{que} pode se \textbf{configurar} de diversas maneiras, inclusive por meio do currículo em ação e de políticas educacionais \textbf{que} envolvam os/as jovens em \textbf{uma} educação integral \textbf{que} promova o seu protagonismo e a tomada de decisões em sua vida acadêmica, principalmente nas cidades ( regiões, territórios, comunidades, e grupos sociais ) mais carentes para \textbf{que} se tornem cidadãos/ãs capazes de colaborar para a construção de \textbf{um} país mais \textbf{justo} e com oportunidades iguais a todos/as.
Conhecer, analisar e avaliar os \textbf{desdobramentos} de tais estudos locais e demais pesquisas em âmbito nacional e internacional combinados aos debates \textbf{que} possam ser suscitados em sala de aula, tanto do/a professor/a em suas múltiplas áreas de conhecimento quanto dos/as estudantes, são práticas fundamentais para se desenvolver \textbf{uma} formação integral.
Refletir sobre o “ território e sua dimensão simbólica ” nas diversas temáticas \textbf{que} podem ser abordadas na escola \textbf{implica} desenvolver \textbf{um} olhar mais crítico sobre os espaços ocupados pelos \textbf{sujeitos}, o processo de formação de tais territórios, as relações sociais estabelecidas na ocupação de tais espaços e os efeitos dessa relação no território em uso.

% matched lemmas: abordagem, agregar, aqui, autor, categoria, conceitual, configurar, conseguir, contudo, crítico-reflexivo, depender, desdobramento, disciplina, dizer, et, exclusão, focar, fronteira, fundamento, implicar, imprescindível, influência, inúmero, justo, lado, lembrar, logo, preciso, predominante, que, revelar, salientar, santo, singularidade, sujeito, sumo, totalidade, um, urgente, ver, viver
\end{textsample}
